\let\negmedspace\undefined
\let\negthickspace\undefined
\documentclass[journal]{IEEEtran}
\usepackage[a4paper, margin=10mm, onecolumn]{geometry}
\usepackage{lmodern}
\usepackage{tfrupee}

\setlength{\headheight}{1cm}
\setlength{\headsep}{0mm}

\usepackage{gvv-book}
\usepackage{gvv}
\usepackage{cite}
\usepackage{amsmath,amssymb,amsfonts,amsthm}
\usepackage{algorithmic}
\usepackage{graphicx}
\usepackage{float}
\usepackage{textcomp}
\usepackage{xcolor}
\usepackage{txfonts}
\usepackage{listings}
\usepackage{enumitem}
\usepackage{mathtools}
\usepackage{gensymb}
\usepackage{comment}
\usepackage[breaklinks=true]{hyperref}
\usepackage{tkz-euclide}
\usepackage{listings}
\def\inputGnumericTable{}
\usepackage[latin1]{inputenc}
\usepackage{color}
\usepackage{array}
\usepackage{longtable}
\usepackage{calc}
\usepackage{multirow}
\usepackage{hhline}
\usepackage{ifthen}
\usepackage{lscape}
\usepackage{tikz}
\usetikzlibrary{patterns}

\begin{document}

\bibliographystyle{IEEEtran}
\vspace{3cm}

\title{4.11.38}
\author{EE25BTECH11061 - V.Sainadh}

\maketitle
{\let\newpage\relax\maketitle}
\renewcommand{\thefigure}{\theenumi}
\renewcommand{\thetable}{\theenumi}
\setlength{\intextsep}{10pt}


{\Large

\textbf{Question:} Find the area of the region bounded by the lines $y=2x+1$, $y=3x+1$ and $x=4$. \\

\textbf{Solution:}  
Given three lines are
\begin{align}
\myvec{-2 & 1}\myvec{x\\y} = 1 \Longrightarrow \vec{n}^\top\vec{x} = 1 \tag{1}
\end{align}
\begin{align}
\myvec{-3 & 1}\myvec{x\\y} = 1 \Longrightarrow \vec{m}^\top\vec{x} = 1 \tag{2}
\end{align}
\begin{align}
\myvec{1 & 0}\myvec{x\\y} = 4 \Longrightarrow \vec{p}^\top\vec{x} = 4 \tag{3}
\end{align}

The three lines form a triangle.  
The vertices of triangle are obtained by 

\textbf{Intersection of : (1) and (3)}
\begin{align}
\myvec{-2 & 1 &\vrule& 1\\
1 & 0 &\vrule& 4} \xrightarrow{R_1 \longrightarrow R_1+2R_2}
\myvec{0 & 1 &\vrule& 9\\
1 & 0 &\vrule& 4}
\end{align}
From the second row we get $y=9$ so $x=4 \Longrightarrow
\vec{A} = \myvec{4\\9}$

\textbf{Intersection of : (2) and (3)}
\begin{align}
\myvec{-3 & 1 &\vrule& 1\\
1 & 0 &\vrule& 4} \xrightarrow{R_1 \longrightarrow R_1+3R_2}
\myvec{0 & 1 &\vrule& 13\\
1 & 0 &\vrule& 4}
\end{align}
From the second row we get $y=13$ so $x=4 \Longrightarrow
\vec{B} = \myvec{4\\13}$

\textbf{Intersection of : (1) and (2)}
\begin{align}
\myvec{-2 & 1 &\vrule& 1\\
-3 & 1 &\vrule& 1} \xrightarrow{R_2 \longrightarrow R_2-R_1}
\myvec{-2 & 1 &\vrule& 1\\
-1 & 0 &\vrule& 0}
\end{align}
From the second row we get $x=0$ so $y=1 \Longrightarrow
\vec{C} = \myvec{0\\1}$

\begin{align}
\vec{A}-\vec{B} = \myvec{0\\-4}, \quad \vec{A}-\vec{C} = \myvec{4\\8}   
\end{align}
\begin{align}
\norm{(\vec{A}-\vec{B})\times(\vec{A}-\vec{C})} 
= \absdet{\myvec{0 & -4\\ 4 & 8}} = |0-(-16)| = 16
\end{align}

\begin{align}
Area \ of \ the \ triangle 
= \frac{1}{2}{\norm{(\vec{A}-\vec{B})\times(\vec{A}-\vec{C})}}
= \frac{16}{2} = 8
\end{align}


Area of the triangle formed by the lines $y=2x+1$, $y=3x+1$ and $x=4$ is $8$ square units.


} % end Large
\newpage 
\textbf{Graph:}
\begin{figure}[h!]
   \centering
   \includegraphics[width=0.8\columnwidth]{figs/Figure_1.png}
   \caption{Representation of a Triangle}
   \label{Plot_1}
\end{figure}


\end{document}
